% Chapter 2: Planning and Learning for Economists
% A narrative introduction to RL theory for PhD economists
% Frames RL as asymptotic approximations of the Bellman fixed-point operator
% Primary references: Bertsekas (2022, 2024), Puterman (1994), Agarwal et al. (2021)

\section{The Theory of Reinforcement Learning}
\label{sec:planning_learning}

\subsection{Introduction}
\label{sec:ch2_intro}

Reinforcement learning and dynamic programming solve the same optimization problem under different information assumptions. Standard treatments of dynamic programming in economics \citep{StokeyLucas1989, Rust1987} assume the transition probabilities and reward functions are known; the computational task is to find the fixed point of the Bellman operator.\footnote{The foundational theory traces to \citet{Bellman1957}, who established the principle of optimality, and \citet{howard1960}, who introduced policy iteration. Applications span consumption-savings problems, job search, investment timing, and structural estimation of discrete choice models.}

Two constraints limit exact dynamic programming. The first is the \emph{Curse of Dimensionality}; the state space $\mathcal{S}$ becomes too large (millions of portfolio configurations, continuous-valued states in DSGE models) to compute and store the value function for every $s \in \mathcal{S}$. The second is the \emph{Curse of Modelling}. The true transition dynamics $P(s'|s,a)$, and often the reward function, are unknown and must be learned from sampled data rather than derived from an analytical prior. Reinforcement learning is not a new theory of decision-making; it is a collection of computational and statistical methods for solving the Bellman equation in the face of these two curses. This chapter develops the theoretical foundations that govern when and why RL algorithms converge.


\subsection{The Geometry of Dynamic Programming}
\label{sec:newton_connection}

Value iteration (VI) and policy iteration (PI) are the workhorses of dynamic programming \citep{StokeyLucas1989}. VI applies the Bellman operator repeatedly until convergence; PI alternates between \emph{policy evaluation} (solving a linear system) and \emph{policy improvement} (taking the greedy action). PI converges faster. Why? The answer reveals a connection between dynamic programming and numerical optimization. Policy iteration is Newton's method applied to the Bellman equation.

Consider the Bellman optimality operator $T$ acting on value functions:
\begin{equation}
(TV)(s) = \max_{a \in \mathcal{A}} \left\{ r(s,a) + \gamma \sum_{s' \in \mathcal{S}} P(s'|s,a) V(s') \right\}.
\label{eq:bellman_operator}
\end{equation}
This operator is nonlinear due to the $\max$. Value iteration applies $T$ repeatedly: $V_{k+1} = TV_k$. Since $T$ is a $\gamma$-contraction in the supremum norm, Banach's fixed-point theorem guarantees $\|V_k - V^*\|_\infty \leq \gamma^k \|V_0 - V^*\|_\infty$. This is Picard iteration, with linear convergence at rate $\gamma$.\footnote{Picard iteration is $x_{k+1} = f(x_k)$ for finding roots of $x = f(x)$. When $f$ is a contraction, convergence is geometric.}

Policy iteration takes a different approach. At the current value estimate $\tilde{V}$, define the greedy policy $\tilde{\pi}(s) = \argmax_a \{ r(s,a) + \gamma \sum_{s'} P(s'|s,a) \tilde{V}(s') \}$. The \emph{policy evaluation} step solves the linear fixed-point equation $V = T^{\tilde{\pi}} V$ exactly, where $T^{\tilde{\pi}}$ is the policy-specific Bellman operator:
\begin{equation}
(T^{\tilde{\pi}} V)(s) = r(s, \tilde{\pi}(s)) + \gamma \sum_{s'} P(s'|s,\tilde{\pi}(s)) V(s').
\end{equation}

The key insight, formalized by \citet{puterman1979}, is geometric: the linear operator $T^{\tilde{\pi}}$ is a supporting hyperplane to the nonlinear operator $T$ at the current iterate\footnote{A supporting hyperplane to a convex function at $x_0$ is a linear function $\ell$ with $\ell(x_0) = f(x_0)$ and $\ell(x) \leq f(x)$ everywhere. The policy operator $T^{\tilde{\pi}}$ plays this role for the Bellman operator.}. Specifically, the operators satisfy tangency ($T^{\tilde{\pi}} \tilde{V} = T\tilde{V}$, agreeing at the current point) and support ($T^{\tilde{\pi}} V \leq TV$ for all $V$, with the linear operator lying below the nonlinear one everywhere). Policy evaluation solves for the fixed point of this linearization exactly. This is precisely the structure of Newton's method; linearize the nonlinear equation at the current point, solve the linearized system, and iterate.\footnote{The Newton interpretation of policy iteration has precursors in \citet{kleinman1968} for Riccati equations in linear-quadratic control and \citet{pollatschek1969} for stochastic games.}\footnote{Algebraically: consider finding the root of $F(V) = V - TV = 0$. For a fixed policy $\pi$, the Bellman operator is affine: $T^\pi V = r^\pi + \gamma P^\pi V$. The Fréchet derivative (Jacobian) of the residual $F$ with respect to $V$ is $(I - \gamma P^{\tilde{\pi}})$, where $\tilde{\pi}$ is the greedy policy at the current iterate. Policy evaluation solves $(I - \gamma P^{\tilde{\pi}})V = r^{\tilde{\pi}}$, which is algebraically the Newton step $V_{k+1} = V_k - [\nabla F(V_k)]^{-1} F(V_k)$. Because the Bellman operator is affine within regions where the greedy policy does not change, Newton's method converges in finitely many steps for tabular MDPs.}

The consequence is quadratic convergence near the optimum. While VI requires $k = \log(100) / \log(1/\gamma)$ iterations to reduce error by a factor of 100, PI typically converges in 5--10 iterations regardless of $\gamma$. At $\gamma = 0.90$, VI needs 44 iterations while PI needs only 5--8; at $\gamma = 0.95$, VI requires 90 versus 5--8; at $\gamma = 0.99$, VI needs 459 iterations while PI still converges in 5--10.

\citet{bertsekas2022newton} extends this interpretation to a broad class of dynamic programming problems. The Newton structure applies whenever the Bellman operator can be written as a pointwise maximum over linear operators: $T = \max_\pi T^\pi$. This includes not only infinite horizon problems with discounting but also optimal stopping problems (job search, option exercise), average cost optimization (inventory, queueing), and minimax formulations for adversarial settings.\footnote{These problems appear under different names in operations research, such as ``stochastic shortest path'' for optimal stopping, ``average-cost MDP'' for long-run average optimization, and ``model predictive control'' (or receding-horizon control) for finite-horizon replanning. See Chapter 1 for historical context on this terminology.} The practical implication is that algorithms with policy-improvement structure (evaluate a policy exactly or approximately, then improve) inherit Newton-like convergence behavior, while pure value-iteration methods (apply $T$ directly) are limited to linear convergence.\footnote{Blackwell's result \citep{blackwell1965} is the fundamental existence theorem for RL. It proves that for discounted MDPs with finite state and action spaces, a stationary deterministic policy $\pi: \mathcal{S} \to \mathcal{A}$ exists that is optimal for all initial states simultaneously. This ``uniform optimality'' has three implications: (1) the search space reduces from history-dependent or stochastic policies to static maps, justifying neural networks that condition only on current state; (2) the optimal policy is independent of the initial distribution $d_0$, so changing where episodes start does not require retraining; (3) PI and VI are guaranteed to converge to the same globally optimal policy regardless of initialization.}


\subsection{Value Learning Methods}
\label{sec:stochastic_approx}

When $P$ is unknown, a single sampled transition $(s, a, r, s')$ can replace the expectation $\mathbb{E}_{s' \sim P(\cdot|s,a)}[V(s')]$. The mathematical foundation is stochastic approximation, developed by \citet{Robbins1952}. Consider the problem of finding $x^*$ such that $g(x^*) = 0$, where $g$ cannot be evaluated directly but one can observe noisy samples $g(x) + \epsilon$. The Robbins-Monro iteration is:
\begin{equation}
x_{t+1} = x_t - \alpha_t [g(x_t) + \epsilon_t],
\label{eq:robbins_monro}
\end{equation}
where $\epsilon_t$ is zero-mean noise. Under two conditions on the step sizes, this converges to $x^*$ with probability one. $\sum_{t=0}^\infty \alpha_t = \infty$ (sufficient exploration, ensuring that learning never ceases) and $\sum_{t=0}^\infty \alpha_t^2 < \infty$ (diminishing noise, ensuring the variance of cumulative updates is finite). The canonical choice $\alpha_t = 1/(t+1)$ satisfies both conditions.

Q-learning \citep{WatkinsDayan1992} is Robbins-Monro applied to the Bellman equation for action-value functions. Define the Q-factor Bellman operator:
\begin{equation}
(FQ)(s,a) = r(s,a) + \gamma \sum_{s'} P(s'|s,a) \max_{a'} Q(s',a').
\end{equation}
The Q-learning update, upon observing transition $(s_t, a_t, r_t, s_{t+1})$, is:
\begin{equation}
Q(s_t, a_t) \leftarrow Q(s_t, a_t) + \alpha_t \left[ r_t + \gamma \max_{a'} Q(s_{t+1}, a') - Q(s_t, a_t) \right].
\label{eq:qlearning}
\end{equation}

The logic of Q-learning is best understood as a Monte Carlo approximation of the Bellman contraction. The true Bellman operator involves an integral over the transition distribution, $(FQ)(s,a) = r + \gamma \int \max_{a'} Q(s',a') \, dP(s'|s,a)$. Since $P$ is unknown, this integral cannot be computed analytically. However, a sample transition $(s,a,r,s')$ acts as a single-point Monte Carlo estimate of this integral. The Q-learning update is simply an exponential moving average (with weight $\alpha_t$) between the current estimate and this noisy Monte Carlo target. Because $F$ is a $\gamma$-contraction in the supremum norm, the expected update drives the estimate toward the fixed point $Q^*$, provided the noise in the Monte Carlo sample averages out over time (which the Robbins-Monro conditions ensure) \citep{tsitsiklis1994}.\footnote{Q-learning can be viewed as root-finding for the expected Bellman residual: the goal is to find parameters such that $\mathbb{E}_{(s,a,r,s')}[\delta_t] = 0$, where $\delta_t = r + \gamma \max_{a'} Q(s',a') - Q(s,a)$ is the temporal difference error. The update is a stochastic gradient step on the mean-squared Bellman error, but with a critical distinction: the gradient is ``semi-gradient'' because the target $\max_{a'} Q(s',a')$ is treated as a fixed constant rather than a function of the parameters being updated. This simplifies computation but disconnects the update from true gradient descent, requiring the specific stability conditions of \citet{tsitsiklis1994}.}

Convergence requires two conditions. Exploration (visiting all state-action pairs infinitely often) ensures identification; every state-action pair must be sampled to learn its value. The Robbins-Monro step-size conditions ($\sum_t \alpha_t = \infty$, $\sum_t \alpha_t^2 < \infty$) balance tracking versus noise suppression. \citet{WatkinsDayan1992} and \citet{jaakkola1994} formalize these; \citet{tsitsiklis1994} provides the general framework. Model-free learning is possible. The optimal value function can be found without ever estimating the transition probabilities.\footnote{Model-free methods are essential when (a) the environment is a physical system with dynamics too complex to write down, or (b) the agent learns directly from interaction. Note that ``model-free'' does not mean ``atheoretical.'' The agent does not store $P(s'|s,a)$ explicitly, but the Q-function serves as an implicit model encoding long-run consequences.}

SARSA provides an on-policy variant. Instead of taking the maximum over next actions, SARSA uses the action actually taken:
\begin{equation}
Q(s_t, a_t) \leftarrow Q(s_t, a_t) + \alpha_t \left[ r_t + \gamma Q(s_{t+1}, a_{t+1}) - Q(s_t, a_t) \right].
\end{equation}
This solves for the value function of the behavior policy $\pi$ rather than the optimal policy. \citet{singh2000} prove convergence under the same step-size conditions, provided the behavior policy converges to a stationary distribution. Q-learning is noisy value iteration on Q-factors; SARSA is noisy \emph{policy evaluation}. The Robbins-Monro conditions ensure that the noise averages out faster than the signal decays, allowing asymptotic convergence despite using only single-sample estimates.

Rollout methods\footnote{A \emph{rollout} simulates one or more trajectories forward from the current state to estimate the value of taking a particular action. The base policy $\mu$ provides a default behavior; the rollout policy improves on $\mu$ by computing one step of lookahead. In practice, this means: from state $s$, try each action $a$, simulate forward using $\mu$ for $H$ steps, average the returns, and pick the action with highest estimated value. AlphaGo's Monte Carlo Tree Search is a sophisticated rollout that grows a search tree rather than simulating independent trajectories.} extend this framework by performing one step of \emph{policy improvement} from a base policy. AlphaZero \citep{Silver2018} implements this principle at scale. Neural networks provide a base value estimate $V^\mu$, and Monte Carlo Tree Search extends the effective lookahead. Given base policy $\mu$, the rollout policy selects:
\begin{equation}
\tilde{\pi}(s) = \argmax_a \left\{ r(s,a) + \gamma \sum_{s'} P(s'|s,a) V^\mu(s') \right\}.
\end{equation}
\citet{bertsekas2021lessons} proves the cost improvement property\footnote{Bertsekas uses cost-minimization notation throughout his work, writing $\min$ where standard RL uses $\max$ and defining value as accumulated cost rather than reward. I translate to the reward-maximization convention used elsewhere in this chapter; the mathematics is equivalent with reversed inequalities.}: $V^{\tilde{\pi}}(s) \geq V^\mu(s)$ for all $s$, with strict inequality unless $\mu$ is already optimal. Each MCTS iteration is a sampled rollout, and the aggregation across many samples approximates the expectation in the Bellman operator.\footnote{\emph{Lookahead} means explicitly considering future states before acting. An $H$-step lookahead builds a tree of depth $H$, evaluating all possible action sequences up to that horizon. The terminal values $V^\mu(s')$ approximate the infinite future beyond the tree. The error bound for $H$-step lookahead with terminal approximation $V^\mu$ is $\|V^{\tilde{\pi}} - V^*\|_\infty \leq \frac{2\gamma^H}{1-\gamma} \|V^\mu - V^*\|_\infty$. At $H=50$ and $\gamma=0.99$, this multiplier is approximately 121, meaning lookahead alone does not compensate for a poor base policy; at $H=500$ the multiplier drops to approximately 1.3, showing that sufficiently deep lookahead can absorb substantial approximation error.}


\subsection{The Central Challenge: The Deadly Triad}
\label{sec:deadly_triad}

The preceding results assume tabular representations, namely, a separate value for each state-action pair. To address the Curse of Dimensionality, one must combine three desirable ingredients, function approximation, bootstrapping, and off-policy learning. Their interaction, however, is the central challenge of reinforcement learning.

When the state space is large or continuous, tabular representations are infeasible. The standard remedy is function approximation. Represent $V(s) \approx \hat{V}(s; \theta)$ using a parameterized family such as polynomials, splines, or neural networks. The parameters $\theta$ are then fitted from sampled Bellman targets. This combination of regression and dynamic programming, however, can diverge---a phenomenon examined below.

Temporal difference (TD) learning with function approximation does not minimize Bellman error directly. Instead, it finds a fixed point of the projected Bellman operator $\Pi T^\pi$, where $\Pi$ projects onto the function approximation subspace. Consider linear function approximation $\hat{V}(s; \theta) = \phi(s)^\top \theta$, where $\phi(s) \in \mathbb{R}^d$ is a feature vector. The projection $\Pi$ maps any value function to its best approximation in the span of the features, under some norm.

The stability of approximate dynamic programming relies on the composition of two operators, the Bellman contraction $T^\pi$ and the projection $\Pi$ onto the function space. Convergence is guaranteed if the combined operator $\Pi T^\pi$ remains a contraction \citep{tsitsiklis1997}. In the on-policy setting, mean-squared error is minimized under the stationary distribution $d^\pi$. Economists can recognize this as OLS. The projection $\Pi$ is orthogonal with respect to the $d^\pi$-norm, defined by $\|V\|_{d^\pi}^2 = \sum_s d^\pi(s) V(s)^2$. Orthogonal projections are non-expansive\footnote{The spectral radius of a linear operator is its largest eigenvalue in absolute value. An operator is non-expansive if its spectral radius is at most 1.} (spectral radius $\leq 1$), so the composition contracts: $\|\Pi T^\pi\| \leq 1 \cdot \gamma < 1$.

In the off-policy setting, samples are drawn from a behavior distribution $\mu \neq d^\pi$. The projection now minimizes error under $\mu$ to approximate a target defined under $d^\pi$. Geometrically, this is an oblique projection, analogous to instrumental variables estimation where the projection and moment conditions use different weighting matrices. Unlike orthogonal projections, oblique projections can expand error (spectral radius $> 1$). If $\|\Pi\|_{d^\pi} > 1/\gamma$, the expansion from projection overwhelms the contraction from discounting. The fixed-point iteration diverges because the operator itself has ceased to be a contraction. \citet{Baird1995} constructs a counterexample: a six-state star MDP with two-dimensional linear function approximation, where samples are drawn uniformly rather than from $d^\pi$. The TD updates cause the weight vector to diverge to infinity. The geometric intuition is that uniform sampling induces a different norm under which $\Pi$ is no longer non-expansive; the composition $\Pi T^\pi$ ceases to be a contraction and can amplify errors.

The divergence phenomenon was recognized before it had a name. \citet{Bertsekas1995_counter} constructed early counterexamples showing that TD($\lambda$) with linear function approximation can diverge even in simple MDPs, establishing that the instability is intrinsic to the combination of bootstrapping and approximation rather than an artifact of particular implementations. \citet{Melo2008} later identified sufficient conditions under which Q-learning with function approximation converges with probability one: the key requirements are that the approximation architecture satisfies a Lipschitz condition\footnote{A function $f$ is Lipschitz if $|f(x) - f(y)| \leq L|x-y|$ for some constant $L$; this bounds how much the function can change.} and that the state-action visitation provides adequate coverage. These conditions are restrictive---they essentially require the function class to be compatible with the MDP structure---but they delineate the boundary between convergent and divergent regimes.

The choice of loss function interacts with the deadly triad in a subtle way. The natural objective is the mean-squared Bellman error (MSBE), $\mathbb{E}_\mu[(Q(s,a) - [TQ](s,a))^2]$, which measures how far the current Q-function is from satisfying the Bellman equation. Computing an unbiased gradient of the MSBE, however, requires two independent next-state samples from the same state-action pair, since the Bellman operator involves an expectation inside a square.\footnote{Expanding $(Q - TQ)^2 = (Q - \mathbb{E}[r + \gamma V(s')])^2$, the gradient involves $\mathbb{E}[r + \gamma V(s')] \cdot \nabla\mathbb{E}[r + \gamma V(s')]$. Estimating this product from a single sample introduces bias because $\mathbb{E}[XY] \neq \mathbb{E}[X]\mathbb{E}[Y]$ in general. A second independent sample from the same $(s,a)$ removes this bias.} This ``double-sampling'' requirement is impractical in most settings. The standard alternative is the mean-squared TD error (MSTDE), which replaces the expected Bellman target with a single-sample estimate. The resulting semi-gradient method targets the projected fixed point $\Pi T^\pi$ rather than the true Bellman fixed point, and it is this mismatch that underlies the instability analyzed above. \citet{Baird1995} proposed residual gradient algorithms that perform true gradient descent on the MSBE, guaranteeing convergence but requiring the double sample. \citet{Antos2008} analyze Bellman-residual minimization in the fitted policy iteration setting, showing that the residual approach achieves near-optimal policies when the function class is rich enough, at the cost of the double-sampling requirement.

The deadly triad has bite because each element is desirable in isolation. Function approximation addresses the curse of dimensionality. When $|\mathcal{S}|$ is astronomical (Go has $10^{170}$ states), storing separate values is infeasible. Bootstrapping enables learning without complete trajectories; TD methods update using current value estimates ($r + \gamma V(s')$) rather than waiting for full returns, essential for continuing tasks and dramatically reducing variance compared to Monte Carlo. Off-policy learning enables Q-learning's key property, namely, learning about the optimal policy while following an exploratory behavior policy, separating data collection from the learning target.

The interaction is fatal. Value-based methods with function approximation no longer apply the true Bellman operator $T$, but a projected operator $\Pi T$. The crux is that $\Pi T$ is not a contraction in any standard norm. When samples come from a distribution $\mu$ that differs from the on-policy distribution $d^\pi$, the projection and contraction are measured in incompatible norms, and errors can compound rather than shrink. Removing any one element restores convergence. Without function approximation (tabular), projection is the identity and Q-learning's contraction applies directly. Without bootstrapping (Monte Carlo returns), targets are independent of current estimates and the problem reduces to supervised learning. Without off-policy learning, samples come from $d^\pi$ and the Tsitsiklis-Van Roy conditions hold.

DQN \citep{mnih2015} addresses the deadly triad through two engineering mechanisms and has since received theoretical justification. Experience replay stores transitions in a buffer and samples mini-batches uniformly, breaking temporal correlation and making the sampling distribution approximately stationary. Target networks hold the target fixed for many updates. Instead of updating toward $r + \gamma \max_{a'} Q(s', a'; \theta)$, DQN updates toward $r + \gamma \max_{a'} Q(s', a'; \theta^-)$, where $\theta^-$ is updated slowly. This decouples the moving target problem from the parameter updates. \citet{zhang2021target} prove that target networks with periodic updates and linear function approximation converge to a neighborhood of the fixed point, with approximation error controlled by the update frequency $\tau$: less frequent updates reduce the neighborhood size at the cost of slower convergence. \citet{Fellows2023} provide a complementary analysis showing that target networks stabilize TD learning by inducing a Wasserstein-type regularization on the value function updates. \citet{Zhang2023} extend these results to the nonlinear setting, proving that DQN with $\varepsilon$-greedy exploration achieves near-optimal sample complexity, closing the gap between Baird's divergence result and DQN's empirical success.

These theoretical developments complement an alternative resolution, gradient TD methods. \citet{sutton2009gtd} reformulate TD as a saddle-point problem\footnote{A saddle-point problem seeks $\min_x \max_y L(x,y)$; gradient TD casts the projected Bellman error as such a problem.}, yielding algorithms (GTD, GTD2, TDC) that perform true stochastic gradient descent on the mean-squared projected Bellman error. These methods converge off-policy with linear function approximation without requiring target networks, but they introduce auxiliary variables and a second set of parameters, increasing per-step computation. In the taxonomy of resolutions, target networks restore stability by slowing the moving target (weakening the bootstrapping component), gradient TD methods restore stability by replacing the semi-gradient with a true gradient (eliminating the projection mismatch), and regularized Q-learning \citep{lim2024regularized} adds an explicit penalty that makes the projected operator contractive. Each approach resolves the triad by modifying a different component of the interaction.


\subsection{When Value-Based RL Works}
\label{sec:when_value_works}

The deadly triad identifies conditions for divergence. The natural question is under what conditions does value-based RL with function approximation succeed? Theoretical advances identify structural assumptions that restore convergence and yield finite-sample guarantees.

The key insight from finite-time analysis of Fitted Value Iteration \citep{MunosSzepesvari2008} is that statistical efficiency requires the sampling distribution $\mu$ to provide adequate coverage of the state-action space visited by near-optimal policies. This is formalized through the concentrability coefficient $C$, which measures the worst-case ratio between the discounted state-action occupancy of any policy and the sampling distribution:
\begin{equation}
C = \max_{\pi} \sup_{s,a} \frac{d^\pi(s,a)}{\mu(s,a)}.
\end{equation}
When $C$ is finite, FVI achieves error bounds scaling polynomially in $C$, the horizon $H = 1/(1-\gamma)$, and the approximation capacity of the function class. The concentrability condition is strong: it requires the offline dataset to cover all states the optimal policy might visit. This is the statistical cost of breaking the contraction principle---the data must compensate for the lost geometric convergence.

The Linear MDP assumption \citep{jin2020} provides an alternative path by imposing structure on the dynamics rather than the data. It posits that transition probabilities and rewards are linear in a known feature map $\phi(s,a) \in \mathbb{R}^d$:
\begin{equation}
P(s'|s,a) = \langle \phi(s,a), \mu(s') \rangle, \quad r(s,a) = \langle \phi(s,a), \theta^* \rangle,
\end{equation}
where $\mu: \mathcal{S} \to \mathbb{R}^d$ is an unknown measure and $\theta^* \in \mathbb{R}^d$ is an unknown parameter vector. Under this assumption, the Bellman operator preserves linearity: applying $T$ to a linear Q-function yields a function that remains linear in the features. This closure property sidesteps the projection error that drives the deadly triad. \citet{jin2020} prove that LSVI-UCB achieves regret $\tilde{O}(\sqrt{d^3 H^3 T})$, polynomial in the feature dimension $d$ and horizon $H$, but independent of $|\mathcal{S}|$ and $|\mathcal{A}|$.

The concept generalizes through Bellman rank \citep{Jiang2017}, which measures the complexity of a value function class under Bellman backups. If the Bellman rank is $M$, then sample complexity scales polynomially in $M$ rather than $|\mathcal{S}|$. Linear MDPs have Bellman rank at most $d$; other structured classes (block MDPs, factored MDPs, linear quadratic regulators)\footnote{Block MDPs group states into latent clusters with shared dynamics; factored MDPs decompose the state into independent components. Both reduce effective dimensionality.} also have low Bellman rank. These results delineate when function approximation is provably safe: polynomial sample complexity is achievable when the environment's dynamics lie within a low-complexity function class.

The contrast with policy gradient methods is instructive. Policy gradients achieve global convergence without concentrability assumptions, but require on-policy data and suffer high variance. Value-based methods can use off-policy data efficiently, but only when the data distribution satisfies coverage conditions or the MDP has low structural complexity. The choice between methods trades off data requirements against computational and statistical efficiency.


\subsection{Policy Learning Methods}
\label{sec:policy_gradient}

Value-based methods find fixed points of the Bellman operator. Policy-based methods take a different approach: parameterize the policy directly as $\pi_\theta(a|s)$ and maximize expected return $J(\theta) = \mathbb{E}_{\pi_\theta}[\sum_{t=0}^\infty \gamma^t R_t]$ by gradient ascent. This formulation sidesteps the Bellman equation entirely and frames reinforcement learning as constrained optimization.

The policy gradient theorem \citep{SuttonMcAllester2000} provides a tractable expression for the gradient:
\begin{equation}
\nabla_\theta J(\theta) = \mathbb{E}_{s \sim d^{\pi_\theta}, a \sim \pi_\theta} \left[ \nabla_\theta \log \pi_\theta(a|s) \, Q^{\pi_\theta}(s,a) \right],
\label{eq:policy_gradient}
\end{equation}
where $d^{\pi_\theta}(s) = (1-\gamma) \sum_{t=0}^\infty \gamma^t \mathbb{P}(s_t = s | \pi_\theta)$ is the discounted state visitation distribution. The fundamental econometric challenge in optimizing $J(\theta)$ is that this distribution depends on $\theta$ through the environment's dynamics. A naive derivative would require $\nabla_\theta d^{\pi_\theta}(s)$, which implies differentiating the unknown transition matrix $P(s'|s,a)$.

The policy gradient theorem sidesteps this entirely via a likelihood ratio (or score function) trick. The theorem transforms a sensitivity analysis problem (how does the system evolve?) into a simpler expectation problem (what is the correlation between the score $\nabla \log \pi$ and the value $Q$?). The gradient can be written as an expectation under the current policy, weighted by action-values, without requiring $\nabla_\theta d^{\pi_\theta}$. The transition dynamics $P(s'|s,a)$ do not appear; the gradient is estimable via sample averages from trajectories alone.

REINFORCE \citep{williams1992} is the simplest policy gradient algorithm. Sample a trajectory $(s_0, a_0, r_0, s_1, \ldots)$, compute the return $G_t = \sum_{k=0}^\infty \gamma^k r_{t+k}$ from each time step, and update:
\begin{equation}
\theta \leftarrow \theta + \alpha \sum_t \nabla_\theta \log \pi_\theta(a_t|s_t) \, G_t.
\end{equation}
This is an unbiased estimator of $\nabla_\theta J(\theta)$, but its variance is high because a single trajectory provides a noisy estimate of $Q^{\pi_\theta}$. Despite high variance, REINFORCE converges to a globally optimal policy in the tabular setting.\footnote{The baseline $b(s)$ subtracted from $G_t$ reduces variance while preserving unbiasedness, since $\mathbb{E}[\nabla_\theta \log \pi_\theta(a|s) b(s)] = 0$ for any baseline independent of $a$.}

Standard gradient descent treats all parameter directions equally. But small changes in $\theta$ can cause large changes in the policy distribution $\pi_\theta$. The natural policy gradient \citep{Kakade2001} accounts for this curvature by preconditioning with the Fisher information matrix:\footnote{The Fisher information $F(\theta) = \mathbb{E}[\nabla \log p \nabla \log p^\top]$ measures curvature of the log-likelihood and appears in the Cramer-Rao bound. Here it measures how policy distributions change with parameters.}
\begin{equation}
\tilde{\nabla}_\theta J(\theta) = F(\theta)^{-1} \nabla_\theta J(\theta), \quad F(\theta) = \mathbb{E}_{s,a} \left[ \nabla_\theta \log \pi_\theta(a|s) \nabla_\theta \log \pi_\theta(a|s)^\top \right].
\end{equation}
Why does NPG recover policy iteration? Standard gradient ascent is sensitive to parameterization: it takes the steepest step in Euclidean parameter space, where units depend on how the policy is parameterized. NPG takes the steepest step in distribution space (measured by KL-divergence), which is invariant to reparameterization. In the tabular case, this geometric correction aligns the gradient exactly with the greedy policy $\tilde{\pi}$ from policy iteration. With step size 1 (and exact estimation), NPG performs one full Newton step; with smaller step sizes, it performs damped Newton updates. This explains its rapid convergence: NPG approximates the quadratic convergence of finding a fixed point rather than the linear convergence of hill-climbing.

The RL objective $J(\theta)$ is non-convex in $\theta$. For economists trained to distrust gradient methods on non-convex objectives, the natural concern is convergence to spurious local optima. The theoretical resolution is striking: for softmax policies, this concern is unfounded. The landscape is ``benign'' in a precise sense. \citet{agarwal2021theory} prove that $J(\theta)$ satisfies a \emph{gradient domination} condition (also called Polyak-\L{}ojasiewicz): whenever $\|\nabla J(\theta)\|$ is small, the policy must be near-optimal. Formally, the sub-optimality $J(\pi^*) - J(\pi_\theta)$ is bounded by a constant times $\|\nabla J(\theta)\|^2$. The implication is immediate: any point where the gradient vanishes is globally optimal. The non-convex landscape has no false peaks, no spurious local maxima. Gradient ascent cannot get trapped.\footnote{However, ``no spurious local optima'' does not imply ``easy optimization.'' The landscape is dominated by vast plateaus (saddle points) where gradients vanish. Without sufficient exploration, the probability of visiting relevant states decays exponentially with the horizon, rendering the gradient exponentially small. Global convergence requires the starting distribution to have adequate coverage relative to the optimal policy's visitation distribution---the ``distribution mismatch coefficient'' of \citet{agarwal2021theory}. The Natural Policy Gradient addresses this by preconditioning with the Fisher Information Matrix, making the update direction covariant: invariant to invertible linear transformations of the parameter space. This standardizes units across parameters, preventing stalling on plateaus caused by poor parameter scaling.}

The Natural Policy Gradient achieves more: \emph{dimension-free} convergence. Standard gradient ascent on $J(\theta)$ has a convergence rate that depends on the smoothness constant, which scales with $|\mathcal{S}|$. NPG circumvents this by preconditioning with the Fisher information matrix $F(\theta)^{-1}$. The key insight is that the state-visitation distribution $d^{\pi_\theta}(s)$ appears in both $\nabla J(\theta)$ and $F(\theta)$; when computing $F^{-1} \nabla J$, these terms cancel analytically. The resulting update rule is equivalent to soft policy iteration and converges at rate $O(1/(1-\gamma)^2 \epsilon)$, independent of $|\mathcal{S}|$ and $|\mathcal{A}|$.

NPG requires computing and inverting the Fisher information matrix $F(\theta)$, which scales as $O(d^2)$ in parameters and is impractical for neural networks. TRPO \citep{Schulman2015} approximates the natural gradient using conjugate gradient methods\footnote{Conjugate gradient is an iterative method for solving linear systems $Ax = b$ without forming $A$ explicitly, requiring only matrix-vector products.} without forming $F$ explicitly, and enforces trust regions via line search.\footnote{Trust region methods build on the performance difference lemma: for any two policies $\pi$ and $\pi'$, $J(\pi') - J(\pi) = \frac{1}{1-\gamma} \mathbb{E}_{s \sim d^{\pi'}} [ \sum_a \pi'(a|s) A^\pi(s,a) ]$, where $A^\pi(s,a) = Q^\pi(s,a) - V^\pi(s)$ is the advantage function. This identity bounds how much policy improvement is possible and motivates constraining updates to regions where advantage estimates remain accurate \citep{Kakade2002}.} PPO \citep{Schulman2017} simplifies further by replacing the hard KL constraint with a clipped surrogate objective, trading theoretical guarantees for computational simplicity.


\subsection{Hybrid Methods}
\label{sec:actor_critic}

REINFORCE estimates policy gradients from sample returns (unbiased, high variance); TD methods use bootstrapped targets $r + \gamma V(s')$ (lower variance, biased when $V$ is approximate). Actor-critic methods combine both. The critic estimates the value function, the actor updates the policy using the critic's estimates.

The theoretical foundation is two-timescale stochastic approximation \citep{konda2000}. Run two concurrent learning processes:
\begin{align}
\text{Critic (fast):} \quad & \theta_{t+1} = \theta_t + \alpha_t^{(c)} \delta_t \nabla_\theta \hat{V}(s_t; \theta_t), \\
\text{Actor (slow):} \quad & \omega_{t+1} = \omega_t + \alpha_t^{(a)} \delta_t \nabla_\omega \log \pi_\omega(a_t|s_t),
\end{align}
where $\delta_t = r_t + \gamma \hat{V}(s_{t+1}; \theta_t) - \hat{V}(s_t; \theta_t)$ is the TD error. The critic updates the value function parameters $\theta$; the actor updates the policy parameters $\omega$.

Convergence requires the critic to learn faster than the actor:
\begin{equation}
\lim_{t \to \infty} \frac{\alpha_t^{(a)}}{\alpha_t^{(c)}} = 0, \quad \text{with both satisfying Robbins-Monro conditions.}
\end{equation}
Under this separation, the actor sees a quasi-stationary critic: from the actor's perspective, the critic provides approximately correct value estimates at each step. The actor's updates are then approximately unbiased policy gradient steps. \citet{konda2000} prove convergence to a local optimum of $J(\omega)$.

A2C (Advantage Actor-Critic) is the synchronous variant: collect a batch of transitions, compute TD errors, and update both networks. A3C \citep{mnih2016a3c} parallelizes this across multiple workers updating a shared parameter server asynchronously. SAC (Soft Actor-Critic) \citep{Haarnoja2018} adds entropy regularization to the actor-critic framework. The agent maximizes the entropy-augmented objective:
\begin{equation}
J_\tau(\theta) = \mathbb{E}_{\pi_\theta}\left[\sum_{t=0}^\infty \gamma^t \left(R_t + \tau \mathcal{H}(\pi_\theta(\cdot|s_t))\right)\right],
\end{equation}
where $\mathcal{H}(\pi) = -\sum_a \pi(a) \log \pi(a)$ is the entropy and $\tau > 0$ is the temperature parameter. The entropy bonus encourages exploration by penalizing deterministic policies. The optimal policy under this objective is softmax in the Q-values: $\pi^*(a|s) \propto \exp(Q^*(s,a)/\tau)$, which matches the logit choice probability from discrete choice econometrics. SAC maintains separate Q-networks (critic) and a policy network (actor), using the soft Bellman operator for critic updates and entropy-augmented policy gradients for actor updates. The actor-critic structure separates identification from optimization. The critic solves a regression problem (estimate $V^\pi$ from data), while the actor solves an optimization problem (improve $\pi$ using the estimated values).


\subsection{Bridging Results}
\label{sec:bridging_results}

Two questions remain: how does approximation error propagate to policy quality, and how does computational complexity scale with problem size? \citet{singh1994} bound the policy degradation from value function errors. If $\hat{V}$ approximates $V^*$ with error $\|\hat{V} - V^*\|_\infty \leq \epsilon$, and $\hat{\pi}$ is the greedy policy with respect to $\hat{V}$, then:
\begin{equation}
\|V^* - V^{\hat{\pi}}\|_\infty \leq \frac{2\gamma}{1-\gamma} \epsilon.
\label{eq:singh_yee}
\end{equation}
At $\gamma = 0.99$, the amplification factor is $2 \cdot 0.99 / 0.01 = 198$. A 1\% error in value function approximation yields at most 198\% error in policy value. This bound is pessimistic but finite: approximate value functions do not cause unbounded policy degradation.

Classical dynamic programming complexity scales with the state space size $|\mathcal{S}|$. For problems like Go, where $|\mathcal{S}| \approx 10^{170}$, exact computation is impossible. \citet{kearns2002} prove that with access to a generative model (a simulator that samples transitions from any state-action pair), near-optimal planning requires samples polynomial in the effective horizon $H = O(\log(1/\epsilon) / \log(1/\gamma))$, the branching factor $|\mathcal{A}|$, and the desired accuracy $1/\epsilon$, with no dependence on $|\mathcal{S}|$.\footnote{The sparse sampling algorithm builds a random tree of depth $H$, estimates values by averaging leaf rewards, and selects the best action.} This explains why MCTS succeeds in large state spaces. Planning complexity depends on horizon and action space, not state space cardinality.


\subsection{Fundamental Tradeoffs}
\label{sec:tradeoffs}

The choice between methods involves distinct tradeoffs rooted in DP structure. Value-based methods target $Q^*$ via the Bellman contraction. In the tabular case convergence is guaranteed, but function approximation introduces the deadly triad. Policy-based methods optimize $\pi_\theta$ directly. Modern theory \citep{agarwal2021theory} establishes global convergence for softmax policies, with high variance as the practical weakness rather than local traps. Actor-critic methods combine both, using the critic for low-variance value estimates while the actor inherits policy gradient's global convergence. Each family traces to DP foundations.

Four additional trade-offs pervade reinforcement learning. First, \emph{exploration versus exploitation}: should the agent act on its current best estimate or gather information to improve future decisions? Naive exploration ($\varepsilon$-greedy) requires samples exponential in the horizon; strategic exploration (UCB, optimism) reduces this to polynomial, formalizing the value of targeted experimentation. Second, \emph{model-based versus model-free}: model-based methods are sample-efficient (each transition updates the entire model) but suffer asymptotic bias if the model class is misspecified; model-free methods are asymptotically unbiased but sample-inefficient, using each transition for a single gradient step. 
Third, \emph{on-policy versus off-policy}: on-policy methods (SARSA, REINFORCE) learn from data generated by the current policy, ensuring stability but discarding past experience; off-policy methods (Q-learning, DQN) reuse stored experience via replay buffers, gaining sample efficiency but risking the instabilities of the deadly triad. Fourth, \emph{bias versus variance in advantage estimation}: REINFORCE uses the full Monte Carlo return (unbiased but high variance); actor-critic methods use bootstrapped TD targets (low variance but biased by the critic's approximation error). The value function baseline is the variance-minimizing choice, motivating the actor-critic architecture as a bias-variance compromise.


\subsection{Conclusion}
\label{sec:ch2_conclusion}

The central insight is that RL algorithms are not mysterious. They are asymptotic approximations to classical dynamic programming operators, justified by the mathematics of contractions, stochastic approximation, and gradient domination. Value iteration becomes Q-learning when expectations are replaced by single samples. Policy iteration becomes the natural policy gradient when the greedy improvement step is approximated by gradient ascent, and NPG recovers PI exactly in the tabular case. The stochastic approximation framework of \citet{Robbins1952} guarantees that under appropriate step-size conditions, noisy iterates converge to the same fixed points as their deterministic counterparts. Reinforcement learning is not a departure from dynamic programming but an extension of it. The central challenge is overcoming the deadly triad. Methods that address it directly, whether through on-policy sampling (PPO), target networks (DQN), or entropy regularization (SAC), have achieved the most significant empirical successes.


